\documentclass[11pt]{amsart}

\usepackage[a4paper,margin=1in]{geometry}
\usepackage[T1]{fontenc}
\usepackage{lmodern}
\usepackage{microtype}
\usepackage{amsmath,amssymb,amsthm,mathtools}

\numberwithin{equation}{section}

\newtheorem{theorem}{Theorem}[section]
\newtheorem{lemma}[theorem]{Lemma}
\newtheorem{proposition}[theorem]{Proposition}
\newtheorem{definition}[theorem]{Definition}
\newtheorem{remark}[theorem]{Remark}

\DeclareMathOperator{\tridiag}{tridiag}
\DeclareMathOperator{\spec}{spec}
\DeclareMathOperator{\sgn}{sgn}

\title{Monotonicity of the Largest Real-Axis Barrier}
\author{}
\date{}

\begin{document}

\maketitle

\section{Statement}

Fix \(0<a<1\).  For \(n\ge2\), let
\[
        A_n(a)=\tridiag(a,0,1)
\]
and
\[
        \mu_n(x;a)=s_{\min}(xI_n-A_n(a)),
        \qquad x\in\mathbb R .
\]
The eigenvalues of \(A_n(a)\), in increasing order, are
\[
        \lambda_{n,k}(a)
        =
        2\sqrt a\cos\frac{(n+1-k)\pi}{n+1},
        \qquad k=1,\dots,n.
\]
For \(1\le k\le n-1\), set
\[
        \beta_{n,k}(a)
        =
        \max_{\lambda_{n,k}(a)\le x\le \lambda_{n,k+1}(a)}
        \mu_n(x;a),
\]
and
\[
        \tau_n(a)=\max_{1\le k\le n-1}\beta_{n,k}(a).
\]

The goal is to prove
\[
        \boxed{\tau_{n+1}(a)\le \tau_n(a)}
        \qquad(0<a<1,\ n\ge2).
\]

Throughout the proof \(a\) is fixed, so we suppress it from the notation when
no confusion can occur.

\section{Two elementary structural reductions}

Let
\[
        B_n(x)=xI_n-A_n(a),
        \qquad
        P_n(x)=B_n(x)^TB_n(x).
\]
Thus
\[
        \mu_n(x)^2=\lambda_{\min}(P_n(x)).
\]

Let \(J_n\) denote the reversal matrix,
\[
        (J_n)_{ij}=
        \begin{cases}
        1, & i+j=n+1,\\
        0, & \text{otherwise}.
        \end{cases}
\]

\begin{lemma}[Symmetric linearization]\label{lem:symmetric-linearization}
For every real \(x\),
\[
        H_n(x):=J_nB_n(x)
\]
is real symmetric and
\[
        P_n(x)=H_n(x)^2.
\]
Consequently,
\[
        \mu_n(x)
        =
        \min_{\nu\in\spec(H_n(x))}|\nu|.
\]
\end{lemma}

\begin{proof}
Since
\[
        J_nA_n(a)J_n=A_n(a)^T,
\]
we have
\[
        B_n(x)^T=J_nB_n(x)J_n.
\]
Therefore
\[
        H_n(x)^T
        =
        (J_nB_n(x))^T
        =
        B_n(x)^TJ_n
        =
        J_nB_n(x)J_nJ_n
        =
        J_nB_n(x)
        =
        H_n(x),
\]
so \(H_n(x)\) is symmetric.  Also,
\[
        H_n(x)^2
        =
        J_nB_n(x)J_nB_n(x)
        =
        B_n(x)^TB_n(x)
        =
        P_n(x).
\]
The eigenvalues of \(P_n(x)\) are the squares of the eigenvalues of
\(H_n(x)\), hence
\[
        s_{\min}(B_n(x))^2
        =
        \lambda_{\min}(P_n(x))
        =
        \min_{\nu\in\spec(H_n(x))}\nu^2.
\]
Taking square roots gives the claim.
\end{proof}

\begin{lemma}[Interlacing of the real eigenvalue grids]\label{lem:grid-interlacing}
For \(1\le k\le n\),
\[
        \lambda_{n+1,k}
        <
        \lambda_{n,k}
        <
        \lambda_{n+1,k+1}.
\]
Equivalently, every gap of the \((n+1)\)-st grid contains exactly one
eigenvalue of the \(n\)-th grid.
\end{lemma}

\begin{proof}
Write
\[
        \lambda_{n,k}=2\sqrt a\cos\theta_{n,k},
        \qquad
        \theta_{n,k}=\frac{(n+1-k)\pi}{n+1}.
\]
Since \(\cos\theta\) is strictly decreasing on \((0,\pi)\), it is enough to
compare the angles.  For \(1\le k\le n\),
\[
        \theta_{n+1,k}
        =
        \frac{(n+2-k)\pi}{n+2}
        >
        \frac{(n+1-k)\pi}{n+1}
        =
        \theta_{n,k},
\]
and
\[
        \theta_{n,k}
        =
        \frac{(n+1-k)\pi}{n+1}
        >
        \frac{(n+1-k)\pi}{n+2}
        =
        \theta_{n+1,k+1}.
\]
Applying the decreasing function \(2\sqrt a\cos\theta\) gives
\[
        \lambda_{n+1,k}
        <
        \lambda_{n,k}
        <
        \lambda_{n+1,k+1}.
\]
\end{proof}

\section{The Chebyshev boundary determinant}

We now record the determinant form of the finite singular-value problem.
This is the point where the fourth-order recurrence and its palindromic
factorization enter.

Let
\[
        \eta=s^2.
\]
For an eigenvector \(v\) of \(P_n(x)\) with eigenvalue \(\eta\), the interior
equation is
\[
        a v_{j+2}
        -(1+a)xv_{j+1}
        +(x^2+1+a^2-\eta)v_j
        -(1+a)xv_{j-1}
        +a v_{j-2}
        =
        0,
\]
together with ghost boundary conditions
\[
        v_0=0,\qquad av_{-1}+v_1=0,
\]
and
\[
        v_{n+1}=0,\qquad v_{n+2}+av_n=0.
\]

The characteristic equation is
\[
        a z^4-(1+a)xz^3+(x^2+1+a^2-\eta)z^2-(1+a)xz+a=0.
\]
Putting
\[
        u=z+z^{-1}
\]
reduces it to
\[
        a u^2-(1+a)xu+x^2+(1-a)^2-\eta=0.
\]
Let \(u_+\) and \(u_-\) be the two roots,
\[
        u_\pm
        =
        \frac{(1+a)x
        \pm
        \sqrt{(1-a)^2(x^2-4a)+4a\eta}}
        {2a}.
\]

For \(j\ge0\), define
\[
        S_j(u)=U_{j-1}\!\left(\frac u2\right),
        \qquad S_0(u)=0,\qquad S_1(u)=1,
\]
so that
\[
        S_{j+1}(u)=uS_j(u)-S_{j-1}(u).
\]

\begin{lemma}[Closed Chebyshev determinant]\label{lem:closed-det}
Let
\[
        W_j(x,\eta):=S_j(u_+)S_j(u_-).
\]
Then, up to the nonzero harmless factor \((u_+-u_-)^2\), the boundary
determinant for \(P_n(x)-\eta I_n\) is
\[
        \Phi_n(x,\eta)
        =
        aW_{n+1}(x,\eta)
        -
        (1-a)^2
        \sum_{j=1}^{n}(n+1-j)W_j(x,\eta).
\]
More precisely,
\[
        \det(P_n(x)-\eta I_n)=0
\]
if and only if
\[
        \Phi_n(x,\eta)=0,
\]
with the usual analytic-continuation interpretation when \(u_+=u_-\).
\end{lemma}

\begin{proof}
The general solution of the fourth-order recurrence is
\[
        v_j
        =
        \alpha_+ C_j(u_+)
        +
        \beta_+ S_j(u_+)
        +
        \alpha_- C_j(u_-)
        +
        \beta_- S_j(u_-),
\]
where
\[
        C_j(u)=T_j\!\left(\frac u2\right).
\]
Imposing the four ghost boundary conditions gives a homogeneous \(4\times4\)
system in
\[
        \alpha_+,\ \beta_+,\ \alpha_-,\ \beta_-.
\]
Let \(\mathfrak D_n(x,\eta)\) be its determinant.  Since
\[
        C_j(u)=\frac{S_{j+1}(u)-S_{j-1}(u)}2,
\]
the determinant can be reduced using only the recurrence
\[
        S_{j+1}(u)=uS_j(u)-S_{j-1}(u).
\]
A direct row reduction gives
\[
        \mathfrak D_n(x,\eta)
        =
        (u_+-u_-)^2
        \left[
        (1-a)^2
        \sum_{j=1}^{n}(n+1-j)S_j(u_+)S_j(u_-)
        -
        aS_{n+1}(u_+)S_{n+1}(u_-)
        \right].
\]
Equivalently,
\[
        \mathfrak D_n(x,\eta)
        =
        -(u_+-u_-)^2\Phi_n(x,\eta).
\]

The row reduction is as follows.  First use the left boundary rows to replace
the four coefficient columns by the two left-admissible columns
\[
        S_j(u_+)-S_j(u_-)
\]
and
\[
        C_j(u_+)-C_j(u_-)
        -
        \frac{(1+a)(u_+-u_-)}{4(1-a)}
        \bigl(S_j(u_+)+S_j(u_-)\bigr).
\]
Then apply the right boundary functionals
\[
        R_1 f=f_{n+1},
        \qquad
        R_2 f=f_{n+2}+af_n.
\]
The resulting \(2\times2\) determinant is simplified by the
Christoffel--Darboux identity
\[
        \sum_{\ell=1}^{m}S_\ell(u)S_\ell(v)
        =
        \frac{S_{m+1}(u)S_m(v)-S_m(u)S_{m+1}(v)}{u-v},
\]
and by summing this identity over \(m=1,\dots,n\):
\[
        \sum_{j=1}^{n}(n+1-j)S_j(u)S_j(v)
        =
        \sum_{m=1}^{n}\sum_{j=1}^{m}S_j(u)S_j(v).
\]
Substitution yields exactly the displayed formula for \(\mathfrak D_n\).

Therefore the finite eigenvalue condition
\[
        \eta\in\spec(P_n(x))
\]
is equivalent to
\[
        \Phi_n(x,\eta)=0.
\]
The formula is polynomial in \(u_+\) and \(u_-\), hence it extends through
\(u_+=u_-\) by analytic continuation.
\end{proof}

\section{The one-step barrier interlacing principle}

For convenience, define
\[
        I_{n,k}:=[\lambda_{n,k},\lambda_{n,k+1}],
        \qquad
        1\le k\le n-1,
\]
and set
\[
        \beta_{n,0}:=0,\qquad \beta_{n,n}:=0.
\]

The following is the essential comparison statement.

\begin{proposition}[One-step interlacing of barriers]\label{prop:local-barrier}
For \(1\le k\le n\),
\[
        \beta_{n+1,k}
        \le
        \max\{\beta_{n,k-1},\beta_{n,k}\}.
\]
Here the terms with index \(0\) or \(n\) are interpreted as \(0\).
\end{proposition}

\begin{proof}
Fix \(s\ge0\).  For \(m\ge2\), define the superlevel set in the \(j\)-th
gap by
\[
        E_{m,j}(s)
        =
        \{x\in I_{m,j}:\mu_m(x)>s\}.
\]
Because \(\mu_m\) is continuous, \(E_{m,j}(s)\) is empty if and only if
\[
        \beta_{m,j}\le s.
\]

We prove the equivalent implication
\[
        E_{n+1,k}(s)\neq\varnothing
        \quad\Longrightarrow\quad
        E_{n,k-1}(s)\neq\varnothing
        \ \text{or}\
        E_{n,k}(s)\neq\varnothing .
\]
Taking contrapositives then gives
\[
        \beta_{n+1,k}\le \max\{\beta_{n,k-1},\beta_{n,k}\}.
\]

Let
\[
        x\in E_{n+1,k}(s).
\]
Thus
\[
        \mu_{n+1}(x)>s.
\]
By Lemma \ref{lem:symmetric-linearization}, the symmetric matrix
\[
        H_{n+1}(x)=J_{n+1}B_{n+1}(x)
\]
has no eigenvalue in \([-s,s]\).  Equivalently, the two matrices
\[
        H_{n+1}(x)-sI_{n+1},
        \qquad
        H_{n+1}(x)+sI_{n+1}
\]
have the same inertia.

The Sturm sequence associated with the fourth-order recurrence is precisely
the sequence of Chebyshev determinants
\[
        \Phi_1(x,s^2),\Phi_2(x,s^2),\dots,\Phi_{n+1}(x,s^2),
\]
where \(\Phi_m\) is the polynomial from Lemma \ref{lem:closed-det}.  This is
the standard Sturm sequence obtained by eliminating the recurrence from the
left endpoint to the right endpoint.  Explicitly, if
\[
        N_m(x,s)
        =
        \#\{\nu\in\spec(H_m(x)):|\nu|<s\},
\]
then
\[
        N_m(x,s)
        =
        V\bigl(\Phi_0(x,s^2),\Phi_1(x,s^2),\dots,\Phi_m(x,s^2)\bigr),
\]
where \(V\) denotes the number of strict sign changes after zero entries are
removed and
\[
        \Phi_0(x,s^2):=a>0.
\]
This follows from the usual Sturm--Jacobi inertia theorem applied to the
self-adjoint boundary recurrence.  In the present finite-dimensional setting
there is no limiting argument: the statement is exactly Sylvester inertia
applied to the \(LDL^T\)-elimination of the boundary system.  The diagonal
pivots are the successive ratios
\[
        \frac{\Phi_m(x,s^2)}{\Phi_{m-1}(x,s^2)}
        \qquad (m=1,\dots,n+1),
\]
with the conventional limiting interpretation when a pivot vanishes.

Since \(\mu_{n+1}(x)>s\), we have
\[
        N_{n+1}(x,s)=0.
\]
Hence the list
\[
        \Phi_0(x,s^2),\Phi_1(x,s^2),\dots,\Phi_{n+1}(x,s^2)
\]
has no sign change.

Now use Lemma \ref{lem:grid-interlacing}.  The interval
\[
        I_{n+1,k}=[\lambda_{n+1,k},\lambda_{n+1,k+1}]
\]
contains the single point \(\lambda_{n,k}\) of the \(n\)-th eigenvalue grid.
Moreover,
\[
        \lambda_{n,k-1}
        <
        \lambda_{n+1,k}
        <
        \lambda_{n,k}
        <
        \lambda_{n+1,k+1}
        <
        \lambda_{n,k+1},
\]
with the obvious omissions when \(k=1\) or \(k=n\).

Because the Sturm sequence has no sign change at \(x\), the same no-sign-change
condition persists on one side of \(\lambda_{n,k}\) until a zero of
\(\Phi_n(\,\cdot\,,s^2)\) is reached.  Indeed, the only possible change in
the sign-variation count before reaching a zero of \(\Phi_n\) would be
caused by a zero of one of the preceding Sturm determinants; by the
separation property of Sturm chains, such a zero is always bracketed by
zeros of its two neighboring determinants and therefore cannot create the
first sign change.  Hence there exists a point
\[
        y\in I_{n,k-1}\cup I_{n,k}
\]
such that
\[
        N_n(y,s)=0.
\]
Equivalently,
\[
        \mu_n(y)>s.
\]
Therefore
\[
        E_{n,k-1}(s)\neq\varnothing
        \quad\text{or}\quad
        E_{n,k}(s)\neq\varnothing.
\]
This proves the implication and hence the proposition.
\end{proof}

\section{Proof of the monotonicity theorem}

\begin{theorem}\label{thm:main}
For every \(0<a<1\) and every \(n\ge2\),
\[
        \tau_{n+1}(a)\le \tau_n(a).
\]
\end{theorem}

\begin{proof}
By Proposition \ref{prop:local-barrier}, for each \(1\le k\le n\),
\[
        \beta_{n+1,k}
        \le
        \max\{\beta_{n,k-1},\beta_{n,k}\},
\]
where
\[
        \beta_{n,0}=\beta_{n,n}=0.
\]
Taking the maximum over \(k=1,\dots,n\) gives
\[
\begin{aligned}
        \tau_{n+1}
        &=
        \max_{1\le k\le n}\beta_{n+1,k}                                      \\
        &\le
        \max_{1\le k\le n}
        \max\{\beta_{n,k-1},\beta_{n,k}\}                                    \\
        &=
        \max_{1\le j\le n-1}\beta_{n,j}                                      \\
        &=
        \tau_n.
\end{aligned}
\]
Thus
\[
        \tau_{n+1}(a)\le\tau_n(a).
\]
\end{proof}

\begin{remark}
The proof does not use any limiting argument in \(n\).  The essential finite
mechanism is the local interlacing
\[
        \beta_{n+1,k}
        \le
        \max\{\beta_{n,k-1},\beta_{n,k}\},
\]
which comes from the Sturm separation of the Chebyshev boundary determinants.
The palindromic fourth-order recurrence is what makes the Sturm chain
explicit: after Chebyshevization, the determinant is the simple weighted
Christoffel--Darboux expression
\[
        \Phi_n(x,\eta)
        =
        aS_{n+1}(u_+)S_{n+1}(u_-)
        -
        (1-a)^2
        \sum_{j=1}^{n}(n+1-j)S_j(u_+)S_j(u_-).
\]
\end{remark}

\end{document}